\section{Formalisation et construction d'une mesure empirique de la charge de travail}
\label{construction-mesure}
\shorthandoff{;:!?}

Construire une mesure n'est pas collecter des données~: c'est trancher une série de décisions de conception, dont chacune engage la validité de tout ce qui suit. L'état de l'art a montré que le TDABC fournit l'équation du coût mais laisse ouvert son intrant principal~: le temps consommé par les activités opérationnelles. Ce chapitre formalise la mesure retenue pour alimenter cette équation dans le contexte du Loan Servicing. Il justifie cinq arbitrages~: retenir l'événement comme unité de charge, séparer l'\textit{Event Processing} et l'\textit{Aftercare}, traiter la séniorité comme une source de biais plutôt que comme une correction de vitesse, conduire une mesure socialement acceptable dans un contexte de nearshoring, et raccorder le temps mesuré à la volumétrie réelle des deals.

La conclusion empirique structurante est volontairement simple~: l'\textit{Event Processing} n'est pas modélisé comme une durée individualisée par opérateur ou par niveau de séniorité, mais comme un temps unitaire fixe. Le temps moyen retenu est

\begin{equation}
\bar{t} \approx 20\ \text{minutes}
\label{eq:temps-unitaire-ep}
\end{equation}

\noindent soit environ 500 observations de clocking. Cette constante devient l'intrant standard du calcul de coût~: qui varie d'un deal à l'autre n'est pas d'abord la vitesse de traitement d'un événement, mais \textbf{le nombre d'événements que le deal génère et la charge d'aftercare associée}.

\subsection{Premier arbitrage~: l'événement comme unité, contre l'équipe et le deal}
\label{sec:arbitrage-unite}

Le choix de l'unité de mesure détermine ce qu'on pourra reconstruire ensuite. Trois granularités étaient possibles, et deux sont des impasses. Mesurer la charge au niveau de l'\textbf{équipe} --- un total d'heures mensuel --- est la donnée que fournissent naturellement les systèmes RH, mais elle est inexploitable pour notre finalité~: on ne peut pas en redescendre au coût d'un deal, puisqu'un deal ne consomme qu'une fraction non identifiable de ce total. Mesurer directement au niveau du \textbf{deal} serait la cible, mais c'est circulaire~: le coût d'un deal est justement l'inconnue~; on ne peut pas le chronométrer globalement sans avoir déjà décomposé ce qui le compose.

Le seul niveau qui rende le problème soluble est l'\textbf{événement opérationnel}. Il est à la fois assez fin pour être chronométrable de façon homogène, et assez structurant pour que tout deal se recompose comme la somme de ses événements. Ce choix aligne la mesure sur le TDABC, dont l'unité est l'activité dotée d'un \textit{duration driver}. Retenir l'événement, c'est faire le pari qu'une charge complexe et continue peut être reconstruite à partir de briques discrètes et mesurables.

L'enquête de clocking confirme que, pour l'\textit{Event Processing}, ce pari est acceptable. Les observations convergent autour d'un ordre de grandeur stable~: un événement de booking consomme en moyenne environ vingt minutes. Les différences fines entre types d'événements existent, mais elles sont moins structurantes, pour le calcul prédictif, que la volumétrie générée par le deal. L'unité pertinente devient donc~: «~un événement vaut environ vingt minutes~», puis le modèle cherche à prédire combien d'événements seront produits.

\subsection{Deuxième arbitrage~: deux objets, non un~--- Event Processing et Aftercare}
\label{sec:arbitrage-decomposition}

La décision la plus lourde de conséquences est de refuser de traiter la charge unitaire comme un bloc homogène. Une lecture naïve mesurerait «~le temps de traitement d'un événement~» et s'arrêterait là. Nos observations imposent de scinder la charge en deux objets qui n'obéissent pas aux mêmes lois~: l'\textit{Event Processing} et l'\textit{Aftercare}.

Le premier, l'\textbf{Event Processing} --- le booking dans Loan~IQ --- correspond à l'exécution opérationnelle d'événements identifiables. Il se prête à une standardisation forte~: sur environ 500 clockings, le temps moyen retenu est une constante empirique $\bar{t}\approx20$ minutes. Le second, l'\textbf{Aftercare}, relève d'une logique différente~: il ne dépend pas seulement de l'existence d'un événement, mais du volume de sollicitations aval que le deal engendre~: courriels, relances, confirmations, réconciliations, interactions avec les prêteurs et reprises opérationnelles. Sa relation aux caractéristiques du deal est quantitative, là où l'Event Processing est traité comme une brique standard multipliée par une volumétrie.
\begin{figure}[htbp]
\centering
\begin{tikzpicture}[
    node distance=1.8cm,
    auto,
    >=stealth',
    stepbox/.style={draw, fill=blue!10, rectangle, rounded corners,
                    minimum width=2.8cm, minimum height=1.3cm,
                    align=center, font=\small},
    epgroup/.style={draw=blue!80!black, dashed, thick, rounded corners,
                    inner sep=0.35cm},
    aftercarebox/.style={draw, fill=green!20, thick, rectangle, rounded corners,
                         align=center, font=\small, text width=10cm}
]

    % --- Étapes du booking (Event Processing) ---
    \node[stepbox] (hobart) {\textbf{1. Hobart}\\ {\footnotesize Création SR / SAR}\\ {\footnotesize (Courriels \& pièces)}};
    \node[stepbox, right=of hobart] (liq) {\textbf{2. Loan IQ}\\ {\footnotesize Booking système}\\ {\footnotesize (Saisie \& validation)}};
    \node[stepbox, right=of liq] (ambre) {\textbf{3. Ambre}\\ {\footnotesize Réconciliation}\\ {\footnotesize (Contrôle final)}};

    % --- Flux séquentiel ---
    \draw[->, thick] (hobart) -- (liq);
    \draw[->, thick] (liq) -- (ambre);

    % --- Encadrement Event Processing ---
    \node[epgroup, fit=(hobart)(liq)(ambre),
          label={[font=\small, align=center]above:\textbf{Event Processing}\\ Temps de booking constant ($\bar{t}\approx 20$ min)}] (ep) {};

    % --- Bloc Aftercare (transversal) ---
    \node[aftercarebox, below=1.4cm of ep] (aftercare)
        {\textbf{Aftercare} — post-traitement unitaire variable\\
         {\footnotesize Échanges de courriels, gestion des prêteurs, réconciliations complexes, interactions Middle Office}};

    % --- Lien EP <-> Aftercare ---
    \draw[<->, thick, dashed] (ep.south) -- (aftercare.north)
          node[midway, right, font=\footnotesize] {Interactions};
\end{tikzpicture}
\caption{Décomposition du traitement d'un événement : Event Processing constant et Aftercare transversal}
\label{fig:flux-ep-aftercare-style}
\end{figure}



Cette dualité est le cœur analytique de notre construction. Confondre les deux objets rendrait la mesure fausse dans les deux sens. Traiter l'Aftercare comme vingt minutes supplémentaires par événement sous-estimerait massivement les gros deals syndiqués, dont la charge aval explose avec le nombre de prêteurs et de facilités. Inversement, chercher à individualiser chaque minute de booking donnerait une précision apparente mais fragile, là où l'évidence empirique permet d'utiliser une constante robuste. La séparation prépare donc directement la modélisation~: prédire le \emph{volume} d'événements pour l'Event Processing, et prédire un \emph{temps} d'Aftercare par régression sur les variables quantitatives du deal.

\subsection{Troisième arbitrage~: écarter la séniorité comme correction de vitesse}
\label{sec:arbitrage-seniorite}

Le principe tayloriste du temps standard invite à chronométrer un temps observé, puis à le corriger pour le ramener à une référence comparable. Une tentation naturelle aurait donc été de remplacer le facteur d'allure subjectif par une variable observable~: la séniorité. L'hypothèse initiale était simple~: un opérateur senior traiterait plus vite qu'un opérateur junior, et il faudrait corriger les temps observés en conséquence.

Cette hypothèse a été abandonnée. La séniorité est trop fortement confondue avec la complexité des dossiers traités. Dans l'organisation observée, les seniors ne prennent pas un échantillon aléatoire de deals~: ils prennent plus souvent les dossiers atypiques, sensibles ou complexes, précisément ceux qui génèrent davantage d'événements, davantage d'exceptions et davantage d'Aftercare. Utiliser la séniorité comme facteur de correction reviendrait donc à attribuer à la vitesse individuelle ce qui relève en réalité de la structure du deal.

Le biais est le suivant~: si un senior met plus longtemps sur un dossier, ce n'est pas nécessairement parce qu'il est moins rapide, mais parce que le dossier qu'il traite est plus complexe. À l'inverse, un junior peut apparaître «~rapide~» parce qu'il traite des opérations plus standardisées. La séniorité ne mesure donc pas proprement une allure~; elle capture simultanément l'expérience, la sélection des dossiers, la criticité opérationnelle et la confiance managériale. C'est un biais de confusion.

Pour cette raison, la séniorité est \textbf{écartée comme correction de vitesse}. Elle peut rester utile descriptivement pour comprendre l'organisation du travail, mais elle ne doit pas intervenir dans la formule empirique de coût. La mesure retient au contraire un temps unitaire commun pour l'Event Processing, $\bar{t}\approx20$ minutes, et laisse la complexité s'exprimer par les variables structurelles du deal~: rôle, prêteurs, facilités, produit, volume d'e-mails et volume d'événements.

\subsection{Quatrième arbitrage~: mesurer un travail dont la mesure menace ceux qui l'exécutent}
\label{sec:arbitrage-social}

Les arbitrages précédents étaient techniques. Celui-ci est social, et il conditionne la validité de tous les autres~: aucune finesse méthodologique ne rachète des temps collectés dans la défiance. Le projet s'inscrit dans une démarche de \textit{cost to serve}, liée à des objectifs d'efficience, de \textit{hard savings} et de \textit{cost avoidance}. Mesurer le coût d'un deal, c'est aussi mesurer la charge des équipes qui le traitent.

Cette dimension est particulièrement sensible dans un contexte de \textit{nearshoring}. Les transferts d'activité entre sites, notamment de Paris vers Lisbonne, ne sont pas une hypothèse abstraite mais une réalité organisationnelle. Le clocking peut donc être perçu comme un instrument de comparaison entre sites, voire comme une justification future de redimensionnement ou de relocalisation. Cette perception n'est pas un bruit extérieur à la mesure~: elle affecte directement ce qui est mesuré.

Le risque principal est l'effet Hawthorne, ou biais de réactivité~: le fait d'être observé modifie le comportement observé. Dans le cas présent, le biais peut aller dans plusieurs directions. Un opérateur peut ralentir parce qu'il se sait observé, accélérer pour donner une image favorable, éviter les cas complexes pendant la période de mesure, ou commenter son travail différemment parce qu'il anticipe l'usage managérial des résultats. La donnée de clocking n'est donc jamais une donnée purement mécanique~; elle est produite dans une relation sociale.

Trois précautions en découlent. Premièrement, dissocier explicitement la mesure de l'évaluation individuelle~: l'objet mesuré est le type d'activité, non la performance nominative. Deuxièmement, expliquer la finalité de la démarche~: rendre visible une charge souvent invisible, notamment l'Aftercare, et non seulement contrôler la productivité. Troisièmement, interpréter les résultats comme des ordres de grandeur robustes plutôt que comme des chronométrages absolus au dixième de minute. Ce cadrage justifie le choix d'une constante empirique de vingt minutes~: elle est plus honnête méthodologiquement qu'une granularité artificielle qui ferait oublier les conditions sociales de production de la donnée.

\subsection{Cinquième arbitrage~: raccorder le temps standard à la volumétrie réelle}
\label{sec:arbitrage-volumetrie}

Un temps standard par événement ne dit rien, seul, du coût d'un deal~: il faut le multiplier par le nombre d'événements que le deal génère. La dernière décision de conception est donc le raccordement de la mesure issue du clocking à la volumétrie issue des systèmes.

L'architecture retenue sépare ce qui est \emph{mesuré une fois} et ce qui est \emph{prédit ou compté en continu}. Le clocking fournit la constante empirique de temps, $\bar{t}\approx20$ minutes. Les systèmes opérationnels fournissent les caractéristiques structurelles des deals et leur historique de volumétrie. La modélisation prédictive exploite ensuite ces caractéristiques pour estimer le volume futur d'événements $\hat{N}(x)$, à partir de variables comme le rôle de la banque, le nombre de prêteurs, le nombre de facilités et le produit.

L'Aftercare est raccordé différemment. Il ne se réduit pas à une multiplication par vingt minutes~: il est estimé comme un temps propre, $\hat{T}_{AC}(x)$, à partir de variables quantitatives~: nombre de prêteurs, nombre de facilités, nombre de mails et nombre d'événements. La mesure produit ainsi deux intrants distincts pour le chapitre de modélisation~: une constante de temps pour l'Event Processing, et une base de variables explicatives pour l'Aftercare.

\subsection{Synthèse~: une mesure empirique robuste et industrialisable}
\label{sec:synthese-mesure}

Les cinq arbitrages se répondent. Le choix de l'événement rend la charge recomposable~; la séparation Event Processing / Aftercare la rend juste~; l'abandon de la séniorité comme correction évite un biais de confusion~; la prise en compte de l'effet Hawthorne et du nearshoring rend explicites les conditions sociales de production de la donnée~; le raccordement à la volumétrie rend la mesure industrialisable.

La mesure obtenue n'est donc pas une table de temps détaillée par opérateur. C'est un modèle empirique plus parcimonieux et plus robuste~: chaque événement d'Event Processing vaut environ vingt minutes, le nombre d'événements est prédit par la structure du deal, et l'Aftercare est estimé séparément. Cette construction fournit l'intrant nécessaire au modèle de coût développé au chapitre suivant.

\shorthandon{;:!?}
